\documentclass[10pt]{article}
\usepackage[margin=20mm]{geometry}
\usepackage{amsmath,amssymb,fontspec,unicode-math,xurl,hyperref}
\IfFontExistsTF{DejaVu Serif}{\setmainfont{DejaVu Serif}}{\setmainfont{Latin Modern Roman}}
\setmathfont{latinmodern-math.otf}
\setlength{\parindent}{0pt}\setlength{\parskip}{6pt}\sloppy
\begin{document}
Abstract: The checked probability identities, loss definitions, trajectory parameterization and benchmark-table arithmetic in GoIRL are coherent. No substantive defect is established in that scope. The correspondence between its fixed-horizon Gibbs model and stationary policy implementation, and two conflicting grid dimensions, remain unresolved. The inspection outcome is Needs Evidence, with no supported Major or Minor finding. Scientific adjudication remains pending.\par
Author: AI-assisted technical review\par
Date:\par
Title: Scoped Technical Verification of GoIRL\par
\# Identity and inspected evidence\par
This report concerns Muleilan Pei, Shaoshuai Shi, Lu Zhang, Peiliang Li and Shaojie Shen, GoIRL: Graph-Oriented Inverse Reinforcement Learning for Multimodal Trajectory Prediction, ICML 2025, PMLR 267, 48699–48713 [Retained latex: \textbackslash{}cite\{paper\}]; assigned 0222/I25-43591, OpenReview `x9LGowGzZu`. The 15-page publisher PDF was acquired on 25 October 2026 UTC. Its SHA-256 is . Page references are one-based PDF pages, matching internal printed numbering.\par
The inspected scope comprises the main method, Algorithm 1, Equations (1)–(17), benchmark definitions and elementary table consistency, and the implementation and evidentiary statements in embedded Appendices A–D. The source manifest records no located/acquired author repository; this does not prove that none was ever released. Possible separate OpenReview material was inaccessible with HTTP 403. No author code, datasets, weights, raw predictions, or supplementary video were executed or inspected. The absence of code is a coverage limit, not itself a scientific defect.\par
\# Checked reasoning\par
\#\# Gibbs likelihood and state visitation\par
For a finite admissible plan set \(\mathcal T\) independent of the reward parameters, with a fixed start/termination convention, write \[p_\theta(\tau)=\frac{\exp R_\theta(\tau)}{Z_\theta},\qquad
 R_\theta(\tau)=\sum_s\mu_\tau(s)R_\theta(s).\] Differentiating the log partition function gives \[\nabla_\theta\overline L_D
 =\sum_s\{\overline\mu_D(s)-\mathbb E_\theta\mu_\tau(s)\}
       \nabla_\theta R_\theta(s).\] This verifies Equation (3), p. 5, provided its expected visitation counts come from the same normalized distribution as Equation (2). Equation (1), p. 4, is the law of total probability over the same plan support.\par
An independently written three-path diagnostic uses count vectors \((2,0),(1,1),(0,2)\), reward parameters \((0.3,-0.2)\), and the first path as a demonstration. It gives analytic gradient \((0.679843332170,-0.679843332170)\). A centered finite difference with step \(10^{-6}\) differs by \(1.51\times10^{-11}\). This checks a finite instance of the identity, not the unpublished driving implementation.\par
\#\# Trajectory and probability formulas\par
Equations (8)–(9), p. 6, use the degree-\(n\) Bernstein basis and \(n+1\) control points. The binomial theorem gives \(\sum_{j=0}^n \binom{n}{j}t^j(1-t)^{n-j}=1\) for \(0\leq t\leq1\); the proposed curve and timestamp substitution are well defined. Equations (15)–(17), p. 13, define the Huber residual loss and pairwise hinge loss consistently with the stated final-displacement winner selection.\par
The product fusion in Equation (11) is a normalized probability vector when both inputs are nonnegative and their product has positive total mass. For inputs \((0.2,0.3,0.5)\) and \((0.5,0.25,0.25)\), the output is \((1/3,1/4,5/12)\). This is a valid heuristic fusion rule; neither Bayesian independence nor a calibration theorem is needed to define it. The report does not infer calibrated probabilities solely from normalization.\par
\#\# Reported numerical relationships\par
The checked combined Argoverse scores equal the stated Brier component plus minFDE: \(0.6232+1.1725=1.7957\) for the single model and \(0.5684+1.1263=1.6947\) for the ensemble. Table 4’s four ablation rows satisfy the same identity at displayed precision. The nuScenes minima and miss rates decrease when more modes are available, as their definitions permit. These are consistency checks of reported numbers, not reproductions.\par
Appendix B and Table 5, p. 13, explicitly list stronger overall ranking scores for QCNet and SEPT and call GoIRL competitive. This qualification is counterevidence against turning the main selected-baseline comparison into a fabricated claim of an undisclosed universal leaderboard lead.\par
\# Unresolved implementation and evidence scope\par
\#\# Finite horizon versus stationary policy\par
Section 3.1, p. 3, defines a fixed planning horizon and calls the grid-cell policy stationary. Algorithm 1, p. 5, returns one policy after \(N\) soft backups; the setup gives \(H=25\), without specifying \(N\), its relation to \(H\), or the complete end-action and terminal-reward conventions.\par
For exact finite-horizon Gibbs sampling, one ordinarily uses remaining-time values and policies, with compatible boundary conditions: \[\begin{aligned}
 V_h(s)&=R(s)+\log\sum_a\exp V_{h-1}(T(s,a)),\\
 \pi_h(a\mid s)&=\exp\{R(s)+V_{h-1}(T(s,a))-V_h(s)\}.\end{aligned}\] A rollout uses decreasing \(h\), or augments the state with remaining time. Reusing a final finite-backup policy throughout a rollout generally prevents the value terms in its path probability from telescoping into the stated Gibbs density.\par
This is Needs Evidence, not a demonstrated erroneous implementation. The algorithm is explicitly called approximate value iteration. A stopping-time model, a converged fixed-point construction with suitable terminal treatment, or implicit time augmentation may explain the description. The required evidence is the actual value/occupancy routine, terminal convention, backup count and sampled plan support. No full algorithm counterexample is asserted in their absence.\par
\#\# Grid specification\par
The implementation paragraph on p. 7 specifies \(100\times100\), downsampled to \(25\times25\). Appendix A.1, p. 12, explicitly specifies \(H_G=W_G=50\) at 1 m spacing for the feature adaptor, separately from a 100 m drivable-area graph. Both statements were visually confirmed. A radius-versus-side interpretation might explain some prose but does not by itself reconcile the two explicit dimensions. The operative crop and downsampling configuration require clarification; no prediction-score consequence has been demonstrated.\par
\#\# Generalization and physical feasibility\par
Section 4.3 and Appendix C provide qualitative drivable-area-change examples. They support the illustrated cases, not a measured success frequency over an independently specified shift population or a universal causal superiority of IRL over behavior cloning. A matched perturbation protocol and quantitative scene-level results would resolve that broader scope.\par
A polynomial proposal is smooth, but unconstrained control points allow unbounded speed or acceleration; learned refinement offsets need not preserve a derivative bound. The paper’s kinematic-feasibility language therefore should not be read as a proved dynamics guarantee. No actual infeasible prediction is established here, and no formal universal theorem was stated.\par
\# Conclusion and state\par
Dimension | Recorded state\par
Checked algebra/arithmetic | No Issue Found within the declared scope\par
Overall inspection outcome | Needs Evidence / scope clarification\par
Supported Major or Minor | None established\par
Scientific adjudication | Pending\par
Repair verification | No repair proposed or executed\par
Author adoption | Unknown; not requested\par
Remaining work | Policy/boundary implementation, grid configuration, raw benchmark and shifted-scene evidence\par
The preserved publisher version is the subject of this review. The checked publisher landing page supplies no correction notice in its returned content; later versions and errata were not exhaustively searched.\par
9 M. Pei et al., GoIRL: Graph-Oriented Inverse Reinforcement Learning for Multimodal Trajectory Prediction, PMLR 267 (2025), 48699–48713. https://proceedings.mlr.press/v267/pei25c.html.\par
\# Reproducibility and actual roles\par
The independent `check.py` and `calc.json` retain the finite-path gradient, Bernstein, fusion and arithmetic diagnostics. Source renders `s05,s07,s12,s13.png` bind the checked formulas and dimensions. The parent performed independent integration, calculations, source-page visual checks, compilation and final-PDF visual inspection. Native auxiliary agents checked theory and empirical/disclosure scope in the same conversation. No human review is asserted. Exact source, evidence, TeX/PDF hashes and compiler/render records reside in the per-paper verification files.\par
\end{document}
